\subsection{Recent Advances in Cybersecurity and Anomaly Detection for Electric Vehicle Supply Equipment}

In the emerging field of EVSE, numerous studies focus either on electric vehicle supply equipment itself or on anomaly detection methods. However, the specific combination of anomaly detection applied to EVSE is not yet widely explored. Therefore, the paper \cite{EV_Security}, titled "Enhancing EV Charging Station Security Using a Multi-dimensional Dataset", is highly relevant to this research. The authors present a dataset used for anomaly detection that includes similar features such as hardware events, power consumption, and other HPC-related data. The paper focuses on various attack scenarios and attempts to detect them using several machine learning models.\\
\\
Another relevant contribution in this field is the paper \cite{paper_relevance_EVSE}, titled "Regression-Based Anomaly Detection in Electric Vehicle State of Charge Fluctuations Through Analysis of EVCI Data" which not only highlights the importance of anomaly detection in EVSE but also provides a comprehensive overview of a realistic charging station scenario. The setup includes four different charging boards, six charging ports, and a battery energy storage system (BESS) handling various types of EVs. The study monitors the state-of-charge (SoC) values of EVs, compares them with expected values, and detects outliers that may indicate anomalies. These detections are performed using machine learning methods within a dynamic environment.\\
\\
In the context of EVSE as part of critical infrastructure, recent research highlights its role as a potential attack vector within complex smart grid environments. Studies report threats such as Denial-of-Service (DoS), Man-in-the-Middle (MitM), spoofing, and data manipulation. Modern approaches increasingly integrate anomaly detection methods using machine learning to identify and mitigate such attacks at an early stage. While not all technical implementation details are relevant here, the findings emphasize the need for advanced security mechanisms to ensure the integrity and availability of EVSE systems \cite{nature2025}.
